\scanpage{52}
% @ lines 1–4 / text
\uline{Proposition} Soient $A$ une ss-Alg. spéciale, et $L$ un sous-fibré inv. de $V$, $\underline L=L^\perp\otimes L\subset M$ le sous-fibré isotrope de $\widetilde{\unsure{\mathfrak g|g}}$ associé. Conditions équivalentes
% @ line 5 / display
\[
 a)\quad \underline L\subset A^\perp
\]
% @ lines 6–7 / text
b) $A$ invarie la droite $L$ (\uline{NB} si $A=\mathcal O_S+J$, il revient au $=$ bien sûr que $J.L\subset L$)
% @ line 7, smaller line below b / text
c) $\exists$ ss-fibré inv. $J$ de $A$ tel que $J.L=0$
% @ line 8 / text
d) $\exists M\subset V$ ss fibré vectoriel \unsure{de rg|droite} tel que
% @ line 9 / display
\[
 \deleted{\unsure{M}}\qquad L^\perp\otimes M\subset A
\]

% @ lines 10–13 / text
De plus, sous ces conditions, le $M$ du c) est unique, et l'unique hom $A\longrightarrow\mathcal O_S$ de noyau $L^\perp\otimes M$ \deleted{\unsure{provient de}} est aussi l'hom $A\xrightarrow{u}\mathcal O_S\simeq\underline{\operatorname{End}}(L)$ déduit de b).
% @ lines 14–17 / text
Enfin, si $J$ est comme dans c), ce noyau est aussi égal à $J$. Et pour $A$ fixé, la corr. $L\mapsto u$ est une bijection entre l'ens des $L$ satisfaisant à a), et l'ens des sections \unsure{d'alg} $A\to\mathcal O_S$

% @ lines 17–19, proof beginning in interline / text
\uline{Dém} La question étant locale, \unsure{sps} $L$ \deleted{\unsure{engendré}} et $L^\perp$ munis de bases $x,x'$, \inserted{\unsure{de} $x'\otimes x$ base de $\underline L$.} si $u$ est \deleted{un} un endom de $V$ i.e section de $M$, on a
% @ line 20 / display
\[
 \operatorname{Tr}u\,x'\otimes x=\langle ux,x'\rangle
\]
% @ lines 21–25 / text
qui est nul ssi $ux$ est orth. à $x'$ i.e. $u(L)\subset L$, d'où a)$\Longleftrightarrow$b). Evidemt si $u:A\to\mathcal O_S\simeq\underline{\operatorname{End}}(\underline L)$ est l'hom correspondant, il est épi donc son noyau $\underline J$ est un ss-fibré inv. tel que $J.\underline L=0$, donc b)$\Rightarrow$c).
% @ lines 26–28 / text
Inv. pour un tel $\underline J$, on sait \inserted{(\unsure{v. leçon.|v. leçons})} que $\mathcal O_S\xrightarrow{\sim}\underline{A/J}$ (car $A/J$ $\mathcal O_S$-Alg. libre de rg $1$) donc $A\simeq J+\mathcal O_S$ et $L$ est stable sous $A$, d'où b)$\Longleftrightarrow$c)
% archive sheet number: 51
